\documentclass[a4paper, 12pt]{ctexart}
\usepackage{graphicx}
\usepackage{geometry}
\usepackage{hyperref}
\usepackage{titlesec}
\usepackage{enumitem}
\usepackage{xcolor}

\geometry{left=2.5cm, right=2.5cm, top=3cm, bottom=3cm}

\title{
    \vspace{-2cm}
    \begin{center}
        \includegraphics[width=5cm]{/home/ubuntu/taistea/frontend/public/logo.png}
    \end{center}
    \vspace{1cm}
    \Huge \textbf{杭州太氏科技有限公司} \\
    \vspace{0.5cm}
    \Large \textbf{设立及运营风险评估与注意事项报告}
}
\author{
    \Large 编制及呈交：太子瑞 \\
    \normalsize (法定代表人、执行董事兼经理，持股 91\%) \\
    \vspace{0.5cm} \\
    \Large 审阅：叶素梅 \\
    \normalsize (股东，持股 9\%)
}
\date{\today}

\begin{document}

\maketitle
\thispagestyle{empty}
\newpage

\tableofcontents
\newpage

\section{公司设立与治理风险}
\subsection{法定代表人与控制权责任}
公司由太子瑞先生担任法定代表人、执行董事兼经理，并持股 91\%，拥有公司的绝对控制权；叶素梅女士（母亲）作为小股东持股 9\%。
\begin{itemize}
    \item \textbf{法定代表人第一责任：} 作为法定代表人与实际控制人，太子瑞不仅享有公司的最高决策权，也承担着公司对外的第一顺位法律责任。若公司发生严重环保事故、食品安全危机、税务偷漏、拒不执行法院判决等极端情况，法定代表人可能面临被限制高消费、行政拘留甚至刑事责任。
    \item \textbf{注册资本与出资义务：} 公司注册资本为 1 万元（或更高认缴额）。根据新《公司法》，股东需在五年内实缴到位。母亲已实缴的金额需通过公对公转账并备注“投资款”，由公司开具《出资证明书》以规避出资瑕疵风险；太子瑞的认缴部分也需规划好出资路径，避免“抽逃出资”的违法行为。
    \item \textbf{财务独立性（公私分明）：} 严禁公司账户与法定代表人（或其母亲）的个人银行卡发生无真实业务背景的频繁资金混同。一旦法院认定“公司人格混同”，有限责任将被击穿，太子瑞将以个人全部财产对公司债务承担连带无限责任。
\end{itemize}

\section{主营业务一：Tais Tea (太氏茶) 与 Fizzwand 销售的合规风险}
\subsection{委托加工（OEM）的核心风险与规避}
公司没有自己的茶叶生产厂和硬件制造厂，必然采取委托加工（代工贴牌）模式。作为品牌方（委托方），公司对最终产品的质量承担不可推卸的首要责任。
\begin{itemize}
    \item \textbf{代工厂资质与责任划分：} 必须严格审查受托生产企业的《营业执照》和《食品生产许可证》（SC证）。在 OEM 合同中，必须加入免责与追偿条款：“因代工厂加工工艺、原料采购等生产环节导致的质量问题、消费者索赔及行政罚款，由代工厂全额承担赔偿责任”。
    \item \textbf{农产品残留红线（GB 2763）：} 茶叶作为入口农产品，农残超标是引发巨额索赔的重灾区。每批次茶叶在装包、出厂、上架销售前，必须要求代工厂提供（或公司自行送检）针对本批次产品的第三方权威机构（CMA/CNAS）出具的农残、重金属及微生物检测合格报告。
    \item \textbf{食品接触材料检测（Fizzwand）：} 搅拌棒作为直接接触饮料的器具，必须由工厂提供符合中国《食品安全国家标准 食品接触材料及制品》（GB 4806 系列）的检测报告。
    \item \textbf{标签与打假人防范：} 包装标签必须严格遵守《预包装食品标签通则》（GB 7718），明确标注委托方与受托方信息。严禁在宣传语中出现“保健”、“降血压”、“防病”等违规词汇，以防被职业打假人依据《食品安全法》索要“退一赔十”。
\end{itemize}

\section{主营业务二：网站一站式制作与数据安全风险}
\subsection{知识产权合规}
\begin{itemize}
    \item \textbf{素材与字体：} 在为客户搭建网站时，使用的图片、图标、音乐、视频、字体和代码组件必须拥有合法的商业使用授权。切勿直接从网络上下载无版权素材，特别是字体（如方正、视觉中国等），极易引发钓鱼诉讼。
    \item \textbf{代码开源协议：} 使用开源代码、框架（如 Vue, React）时，需严格遵守其开源协议（MIT, Apache 等），避免违反传染性协议（如 GPL）导致客户代码被迫开源。
\end{itemize}
\subsection{数据安全与网络合规}
\begin{itemize}
    \item \textbf{ICP 备案合规：} 公司自行运营的商城网站及为国内客户制作并托管的网站，均需要进行 ICP 备案。涉及电商交易的，还需办理 EDI（增值电信业务经营许可证）或进行公安联网备案。
    \item \textbf{个人信息保护：} 网站需制定合规的《隐私政策》和《用户协议》，在收集用户的姓名、电话、收货地址、支付信息时必须获取用户的明确同意，并做好数据加密保护，防止数据泄露。
\end{itemize}

\section{进出口与跨境电商业务风险}
如果 Tais Tea 或 Fizzwand 涉及直接的海外在线销售或需要从海外采购原材料：
\begin{itemize}
    \item \textbf{合法付汇与资金出海：} 公司对外的任何支付（如购买海外 AWS 服务器、支付 Google 广告费）必须通过真实的“服务贸易”或“货物贸易”合同，在银行正规渠道完成付汇。严禁通过“地下钱庄”或买卖加密货币（USDT）进行资金跨境转移，以免触犯“洗钱”或“非法经营罪”的刑法红线。
    \item \textbf{进出口资质与海关备案：} 开展实体货物进出口，公司需前往商务局、海关、外汇管理局办理进出口收发货人备案，获取外贸经营权。
    \item \textbf{目的国严苛准入标准：} 
        \begin{itemize}
            \item \textbf{欧盟：} 对茶叶的农残要求（EU MRLs）为全球最严，诸多国内合格的茶叶无法进入欧盟市场；
            \item \textbf{美国：} 涉及食品和食品接触器具（如 Fizzwand）出口至美国，需进行 FDA 设施注册并在发货前进行事先申报（Prior Notice）。
        \end{itemize}
    \item \textbf{关税与物流风险：} 跨境电商需充分了解目的国的关税起征点和增值税（VAT/GST）政策，避免货物在海关被扣留或产生高昂的退运成本。
\end{itemize}

\section{税务合规与财务风险}
\begin{itemize}
    \item \textbf{在线销售的“无票收入”入账：} 在电商平台（如淘宝、微信小程序或自建商城）的销售收入，无论买家是否索要发票，公司均必须如实申报纳税。在“金税四期”大数据监管下，第三方支付流水与公户绑定的情况下，隐瞒无票收入极易被税务稽查认定为偷漏税。
    \item \textbf{合理享受小微企业红利：} 目前国家对月销售额 10 万元以下的增值税小规模纳税人免征增值税；对小型微利企业也有相应的企业所得税减免（如针对前 100 万利润有优惠税率）。应聘请专业代账机构进行合规的税收筹划。
    \item \textbf{成本费用的发票收集：} 委托工厂加工、购买服务器、甚至是合规的办公消耗，均需向对方索要合规的增值税发票，以便在企业所得税前抵扣。无票支出将导致公司账面利润虚高，从而多缴所得税。
\end{itemize}

\section{总结}
作为杭州太氏科技有限公司的法定代表人兼绝对控股股东，太子瑞先生处于公司权力与责任的顶端。在冲刺业绩的同时，务必筑牢合规防火墙：严守资金公私分离底线，把控代工产品的农残与食品安全生命线，规范税务申报。建议在开展核心业务前，与代工厂签订权责清晰的法务合同，并委托专业代账公司处理财税申报，确保公司在合法合规的轨道上稳健发展。

\end{document}